\documentclass[10pt,a4paper]{amsart}
\usepackage[margin=19mm]{geometry}
\usepackage{amsmath,amssymb,mathtools}
\usepackage[scr=rsfs,cal=euler]{mathalfa}
\usepackage{xfrac}
\usepackage[unicode,hidelinks,bookmarks=false]{hyperref}
\newcommand{\ie}{i.e.\ }
\newcommand{\cf}{cf.\ }
\newcommand{\resp}{respectively\ }
\newcommand{\Subsection}{\subsection}
\providecommand{\nobreakdash}{\nobreak\hbox{-}}
\setlength{\emergencystretch}{2em}
\multlinegap0pt
\usepackage{amssymb}
\usepackage{bm}
\usepackage{tikz}
\usepackage{float}
\newtheorem{theorem}{Theorem}
\newcommand{\A}{\mathbb A}
\newcommand{\C}{{\mathbb C}}
\newcommand{\F}{{\mathbb F}}
\newcommand{\Spec}{\mbox{Spec}}
\newcommand{\Z}{{\mathbb Z}}
\title{Remarks on polynomial count varieties}
\author{Nicholas M. Katz, Fernando Rodriguez Villegas}
\date{Source: arXiv:2603.07062v1 (CC BY 4.0)}
\begin{document}
\maketitle

\noindent\emph{Source and licence.} Remarks on polynomial count varieties. Version arXiv:2603.07062v1 (CC BY 4.0), distributed by arXiv under the Creative Commons Attribution 4.0 International licence, http://creativecommons.org/licenses/by/4.0/, as recorded in the arXiv licence metadata for this exact version and shown on the abstract page https://arxiv.org/abs/2603.07062v1. Full licence notice: \texttt{licenses/arxiv-2603.07062-CC-BY-4.0.txt}.\par\smallskip

\noindent\textit{Editorial scope.} Counted unit: the article's theorem simpl-count-thm (source0.tex lines 451--489). The article's complete original proof of it is delivered with it (source0.tex lines 490--522, one \texttt{proof} environment of 33 lines). No mathematics is written, repaired or re-derived: the statement and the proof are the article's own lines, in the article's own notation, and no retained line is substituted or re-flowed.  Retained uncounted as the source-local context the proof needs: lines 433--436.  Nothing was omitted from any retained span, so the package carries no omission record.  No retained line contains a citation, so the wrapper carries no bibliography and no bibliography entry.  No retained line contains a figure environment, an image-inclusion command or drawing code, so no image is delivered and the package declares no asset dependency.  Editorial formatting only, in addition: an AMS \texttt{amsart} preamble that restates the article's own declarations and macros used by the retained lines, namely the environments \texttt{theorem} on separate counters, together with the standard packages those lines need.  The retained environments print in this document as Theorem 1 in order of appearance, whereas the article prints the same items with the numbering of its own full document.  The complete native source of the article as distributed in the arXiv e-print for this exact version is bundled unchanged as \texttt{sources/195914/source0.tex}.
\begin{equation}
\label{f-defn}
f_{r,n}:=\#\{S \,|\, \#S=r, \dim \Delta_S=n-1\},
\end{equation}

\begin{theorem}
\label{simpl-count-thm}
  Let $X\subseteq \A^N$ be the zero locus over $\C$ of a polynomial
  $H\in \Z[x_1,\ldots,x_N]$ with $N\geq 1$ whose Netwon polytope $\Delta$
  is equivalent to $\sigma_N$. Let
$$
X':=X\cap T,
$$
where 
$$
T:=\Spec(\Z[x_1,x_1^{-1},\ldots,x_N,x_N^{-1}])
$$
is the torus. 

Then the varieties $X'$ and  $X$ are polynomial count outside $D$, for
$$D:={\rm\ the \ product\ of \ all\  nonzero\  coefficients\ of\ }H.$$
More
precisely, 

\begin{equation}
\label{simpl-count-0}
\#X'(\F_q)= c_N(q), 
\end{equation}
and
\begin{equation}
\label{simpl-count}
\#X(\F_q)=\sum_{r=0}^N\sum_{n=0}^N
f_{r,n}c_n(q)(q-1)^{r-n},
\end{equation}
where
$$
c_n(q)=
\begin{cases}
\sum_{i=1}^n(-1)^{n-i}\binom{n}{i}\, q^i + (-1)^n&\quad n\geq 0\\
1 & \quad n=0
\end{cases},
$$
 and the $f_{r,n}$ are defined in~\eqref{f-defn}.
\end{theorem}

\begin{proof}
  It is easy to verify that for $N\ge 1$
$$
\#\{(x_1,\ldots,x_N)\in (\F_q^\times)^N \,|\,
x_1+\cdots
+x_N=0\}= c_N(q)=\sum_{n=1}^N(-1)^{N-n}\binom Nn\, q^{n-1} + (-1)^N
$$
by M\"obius inversion as
$$
\sum_{n=0}^N\binom N nc_n(q)=q^{N-1},\qquad N\geq 1.
$$
Since $\Delta$ is equivalent to the standard
simplex~\eqref{simpl-count-0} follows. Indeed, 
let $y_i:=x^{m_i}$, where $m_1,\ldots,m_N$ are
the vertices of $\Delta$. Given the equivalence of $\sigma_N$ and
$\Delta$ this yields an automorphism $(x_1,\ldots,x_N)\mapsto
(y_1,\ldots,y_n)$ of the torus. But then the two equations
$$
x^{m_1}+\cdots+x^{m_N}=0, \qquad \qquad y_1+\cdots+y_N=0
$$
have the same number of solutions in $\F_q^\times$.

To prove~\eqref{simpl-count} we proceed recursively. Notice that
specializing $H(x_1,\ldots,x_N)$ by setting $x_i=0$ for $i\notin S$
with $S\subseteq \{1,\ldots,N\}$ results in a polynomial
$H_S(x_1,\ldots,x_N)$ whose Newton polynomial $\Delta_S$ is a face of
$\Delta$ (or empty). Therefore, $\Delta_S$ is also equivalent to
$\sigma_n$ for some $n\geq 0$ or empty (if $H_S$ vanishes
completely). In particular, if $\Delta_S$ is non-empty the number of
solutions of $H_S(x_1,\ldots,x_N)$ with $x_i\in \F_q^\times$ is
$c_n(q)(q-1)^{r-n}$, where $r:=\#S$. On the other hand, if $\Delta_S$
is empty then the number of solutions is $(q-1)^r$. The claim follows.
\end{proof}

\end{document}
